\section{The \sys{} Framework and the World-Action Design Space}
\label{sec:framework}

This section defines the object of study. \sref{sec:jointfm} writes a
world-action model as one joint flow-matching problem over a video stream and an
action stream and derives the property that makes the deployment-time sampler a
free variable. \sref{sec:coupling} enumerates the ways the action stream can be
attached to a video backbone. \sref{sec:masks} formalises cross-modal visibility
as a block mask and derives how that choice interacts with action-free
pretraining data. \sref{sec:repr} specifies the visual representation interface
and the reducer that makes semantic and reconstructive latents comparable.
\sref{sec:contract} gives the observation and action contract that lets one head
serve many embodiments. \sref{sec:deploy} derives the inference-time schedules
that the framework exposes at deployment.

\subsection{A World-Action Model as One Joint Flow-Matching Problem}
\label{sec:jointfm}

\Paragraph{Inputs and targets.}
A training window carries a clip of $N_v$ observation frames, a language
instruction $c$, a proprioceptive state $s$, and an action chunk
$a_0 \in \mathbb{R}^{H \times d}$ of horizon $H$. A frozen visual encoder
$E$ maps the clip to a latent grid
$z_0 \in \mathbb{R}^{C \times T \times H' \times W'}$. The first frame is
encoded separately and pinned into the latent grid as a clean conditioning
prefix, so the model always predicts a future conditioned on a present
observation.

\Paragraph{Per-stream flow matching.}
Both streams use the same rectified-flow construction
\citep{lipman2023flowmatching,liu2023rectifiedflow}. For a clean sample $x_0$,
Gaussian noise $\epsilon \sim \mathcal{N}(0, I)$ and a noise level
$\sigma \in [0,1]$,
\begin{equation}
  x_\sigma \;=\; (1-\sigma)\, x_0 \;+\; \sigma\, \epsilon,
  \qquad
  u(x_0, \epsilon) \;=\; \epsilon - x_0 ,
  \label{eq:fm}
\end{equation}
where $u$ is the constant velocity of the straight probability path joining the
data point to the noise sample. Noise levels are drawn from a shifted grid. Let
\begin{equation}
  f_\lambda(t) \;=\; \frac{\lambda t}{1 + (\lambda - 1) t},
  \qquad \lambda > 0, \; t \in [0,1] ,
  \label{eq:shift}
\end{equation}
denote the time reparameterisation that is informationally equivalent to
scaling the latent variance by $\lambda$ \citep{esser2024sd3}. The grid of a
stream with $K$ levels and shift $\lambda$ is
$\sigma_k = f_\lambda\!\left(1 - k/K\right)$ for $k = 0, \dots, K-1$, and
training uses $K = 1000$ with $\lambda_V = \lambda_A = 5$.

\Paragraph{The joint objective.}
One network $F_\theta$ consumes both noisy streams together with the
conditioning prefix, the instruction and the state, and returns two velocity
fields, as \fref{fig:overview} shows,
\begin{equation}
  \left(\hat{v}, \hat{a}\right)
  \;=\;
  F_\theta\!\left(z_{\sigma_v},\, a_{\sigma_a},\, z^{\text{cond}},\, c,\, s,\, \sigma_v,\, \sigma_a\right).
  \label{eq:net}
\end{equation}
Writing $\|r\|^2_{M} = \sum M \odot r^2 \big/ \sum M$ for the mask-normalised
squared error of a residual $r$, the training loss is
\begin{equation}
  \mathcal{L}(\theta)
  \;=\;
  \lambda_V \, \mathbb{E}\Big[\, w_V(\sigma_v)\,
    \big\| \hat{v} - u(z_0, \epsilon_V) \big\|^2_{M_V} \Big]
  \;+\;
  \lambda_A \, \mathbb{E}\Big[\, w_A(\sigma_a)\,
    \big\| \hat{a} - u(a_0, \epsilon_A) \big\|^2_{M_A} \Big].
  \label{eq:loss}
\end{equation}
The video mask $M_V$ excludes the clean conditioning latents and any padded
frame. The action mask $M_A$ is a per-element mask over horizon steps and action
dimensions, which is the mechanism that lets one action head absorb many
embodiments (\sref{sec:contract}). The per-level weights follow the
bell-shaped profile
\begin{equation}
  w(\sigma) \;\propto\; \exp\!\left(-2\left(\sigma - \tfrac{1}{2}\right)^{2}\right),
  \label{eq:weight}
\end{equation}
normalised to unit mean over the grid, so intermediate noise levels dominate the
gradient and the two streams share a consistent level weighting.

\Paragraph{The two noise levels are independent, which frees the sampler.}
The single design decision that separates \eqref{eq:loss} from a conventional
diffusion policy is that $\sigma_v$ and $\sigma_a$ are drawn independently.

\begin{proposition}[Product coverage]
\label{prop:product}
Let $\sigma_v \sim \mathrm{Unif}\{\sigma^V_0, \dots, \sigma^V_{K-1}\}$ and
$\sigma_a \sim \mathrm{Unif}\{\sigma^A_0, \dots, \sigma^A_{K-1}\}$ be sampled
independently for every training example. Then the empirical distribution of
level pairs converges to the uniform measure on the product grid
$\{\sigma^V_k\} \times \{\sigma^A_j\}$, and the minimiser of \eqref{eq:loss}
matches the conditional expectations
$\mathbb{E}[u(z_0,\epsilon_V) \mid z_{\sigma_v}, a_{\sigma_a}]$ and
$\mathbb{E}[u(a_0,\epsilon_A) \mid z_{\sigma_v}, a_{\sigma_a}]$ at every pair on
that grid.
\end{proposition}

\begin{proof}[Proof sketch]
Independence makes the joint level distribution the product of the two marginals,
so every pair receives probability $1/K^2$ and appears infinitely often in the
limit. Conditioned on a fixed pair, \eqref{eq:loss} reduces to two weighted
least-squares regressions whose minimisers are the stated conditional
expectations. Since the losses at distinct pairs are additive and the pair is an
input to $F_\theta$, the pairwise minimisers are attained simultaneously.
\end{proof}

Every monotone path from $(\sigma_v, \sigma_a) = (1,1)$ to $(0,0)$ inside the
unit square is therefore an admissible sampler, because each of its points is a
level pair the model was trained on. Tying the two levels together during
training would supervise only the diagonal of the square and would fix the
sampler at training time. \sref{sec:deploy} turns this freedom into a
parameterised family of trajectories, and \sref{sec:q6} measures it.

\begin{figure}[t]
\centering
\adjustbox{max width=\linewidth}{%
\begin{tikzpicture}[font=\small,
  box/.style={draw=black!55, rounded corners=2pt, align=center, inner sep=3pt, minimum height=8mm},
  enc/.style={box, fill=tianzheblue!12},
  core/.style={box, fill=tianzhegreen!14, minimum height=22mm},
  head/.style={box, fill=tianzhered!12},
  ar/.style={-{Stealth[length=2mm]}, draw=black!65}]

% inputs
\node[enc] (frames) at (0,1.55) {observation clip\\\scriptsize $N_v$ frames};
\node[enc] (first)  at (0,0.45) {present frame};
\node[enc] (lang)   at (0,-0.65) {instruction $c$};
\node[enc] (prop)   at (0,-1.75) {state $s$};

% encoders
\node[enc] (venc) at (3.3,1.0) {frozen visual\\encoder $E$};
\node[enc] (tenc) at (3.3,-1.2) {frozen text\\encoder};

\draw[ar] (frames.east) -- (venc.west);
\draw[ar] (first.east) -- (venc.west);
\draw[ar] (lang.east) -- (tenc.west);

% noisy streams
\node[box, fill=tianzheblue!8] (zs) at (6.7,1.8) {$z_{\sigma_v}$\\\scriptsize noisy video latents};
\node[box, fill=tianzheorange!14] (as) at (6.7,0.35) {$a_{\sigma_a}$\\\scriptsize noisy action chunk};
\node[box, fill=black!5] (ctx) at (6.7,-1.2) {context tokens\\\scriptsize text and state};

\draw[ar] (venc.east) -- (zs.west);
\draw[ar] (tenc.east) -- (ctx.west);
\draw[ar] (prop.east) -- (5.05,-1.75) |- (ctx.west);

% core
\node[core] (core) at (10.3,0.45) {coupled\\denoiser $F_\theta$\\\scriptsize \sref{sec:coupling}, \sref{sec:masks}};
\draw[ar] (zs.east) -- (core.west);
\draw[ar] (as.east) -- (core.west);
\draw[ar] (ctx.east) -- (core.west);

% outputs
\node[head] (vhat) at (13.9,1.35) {video velocity $\hat{v}$};
\node[head] (ahat) at (13.9,-0.45) {action velocity $\hat{a}$};
\draw[ar] (core.east) -- (vhat.west);
\draw[ar] (core.east) -- (ahat.west);

% level annotations
\node[anchor=west, font=\scriptsize, text=tianzhebrown] at (5.40,2.85)
  {$\sigma_v$ drawn from the shifted grid};
\node[anchor=west, font=\scriptsize, text=tianzhebrown] at (5.40,-2.25)
  {$\sigma_a$ drawn from the shifted grid, independently of $\sigma_v$ (Prop.~\ref{prop:product})};
\end{tikzpicture}}
\caption{One forward pass of a world-action model in \sys{}. Frozen encoders
map the observation clip and the instruction into a latent grid and a context
sequence. Two streams are noised at independently drawn levels and denoised
together, and the network returns one velocity field per stream. Independent
level sampling is what makes the deployment-time trajectory through the level
square a free variable.}
\label{fig:overview}
\end{figure}

\Paragraph{Sampling.}
At deployment the policy walks a schedule
$\mathcal{S} = \{(\sigma_v^{(k)}, \sigma_a^{(k)})\}_{k=0}^{K'}$ that ends at
$(0,0)$, taking one Euler step per stream and per iteration,
\begin{equation}
  x^{(k+1)} \;=\; x^{(k)} \;+\; \big(\sigma^{(k+1)} - \sigma^{(k)}\big)\,
  \hat{u}^{(k)} .
  \label{eq:euler}
\end{equation}
\begin{algorithm}[t]
\caption{Joint denoising with a free level schedule and a velocity cache}
\label{alg:sample}
\begin{algorithmic}[1]
\STATE \textbf{input:} present frame, instruction $c$, state $s$, schedule
  $\mathcal{S}$, visibility mask $M$, cache threshold $\tau$, skip budget $m$
\STATE $z \leftarrow$ Gaussian noise; pin $z^{\text{cond}}$ into the conditioning prefix of $z$
\STATE $a \leftarrow$ Gaussian noise; $\; j \leftarrow 0$
\FOR{$k = 0$ \TO $K' - 1$}
  \STATE $(\sigma_v, \sigma_a) \leftarrow \mathcal{S}[k]$;\enspace
         $(\sigma_v', \sigma_a') \leftarrow \mathcal{S}[k+1]$
  \IF{$k \geq 2$ \AND $j < m$ \AND
      $\cos\big(\hat{v}^{(k-1)}, \hat{v}^{(k-2)}\big) > \tau$}
    \STATE $\hat{v}^{(k)} \leftarrow \hat{v}^{(k-1)}$;\enspace
           run the action branch alone to obtain $\hat{a}^{(k)}$;\enspace
           $j \leftarrow j + 1$
  \ELSE
    \STATE $\big(\hat{v}^{(k)}, \hat{a}^{(k)}\big) \leftarrow
            F_\theta\big(z, a, z^{\text{cond}}, c, s, \sigma_v, \sigma_a; M\big)$;\enspace
           $j \leftarrow 0$
  \ENDIF
  \STATE $z \leftarrow z + (\sigma_v' - \sigma_v)\, \hat{v}^{(k)}$;\enspace
         restore $z^{\text{cond}}$ in the conditioning prefix
  \STATE $a \leftarrow a + (\sigma_a' - \sigma_a)\, \hat{a}^{(k)}$
\ENDFOR
\STATE \textbf{return} the action chunk $a$ in physical units after inverse normalisation
\end{algorithmic}
\end{algorithm}
Video decoding is optional, since a control loop consumes only the action
stream. \sref{sec:deploy} specifies the schedule family and the cache.

\subsection{How the Action Stream Attaches: Six Coupling Architectures}
\label{sec:coupling}

A video backbone is a stack of $L$ transformer blocks. Block $\ell$ factors into
a pre-attention map that produces queries, keys and values, a self-attention, and
a post-attention map that applies the output projection, the feed-forward network
and the residual updates. \sys{} exposes this factorisation, so an action stream
may be inserted at any of the three points. Six concrete insertions follow, and
\fref{fig:families} draws them.

\Paragraph{Single system.}
The action chunk is embedded into $S_A = H$ tokens of the backbone width $d_V$,
the proprioceptive state is embedded into one token, and both are concatenated
onto the video token sequence. The shared blocks process the concatenation and a
small multilayer perceptron decodes the trailing action segment. The variant
\texttt{vanilla} keeps the shared capacity as it stands. The variant \texttt{moe} adds a
modality expert at selected layers,
\begin{equation}
  h_A^{(\ell)} \;\leftarrow\; h_A^{(\ell)} \;+\; \mathrm{FFN}^{\text{exp}}_\ell\!\left(h_A^{(\ell)}\right),
  \qquad \ell \in \mathcal{B},
  \label{eq:moe}
\end{equation}
where $h_A^{(\ell)}$ is the action segment of the layer output and $\mathcal{B}$
is the selected layer set. The single system reuses every video parameter for the
action computation, which makes it the cheapest coupling and ties the two
modalities to one hidden width.

\Paragraph{Dual system with joint self-attention.}
Two experts run side by side: a video expert of width $d_V$ and an action expert
of width $d_A$. At every layer both experts project their own residual streams
into a shared attention space of $n$ heads of width $d_h$, the two token
sequences are concatenated, one attention runs over the concatenation, and the
result is split back,
\begin{equation}
  \begin{bmatrix} O_V \\ O_A \end{bmatrix}
  \;=\;
  \mathrm{softmax}\!\left(
    \frac{1}{\sqrt{d_h}}
    \begin{bmatrix} Q_V \\ Q_A \end{bmatrix}
    \begin{bmatrix} K_V \\ K_A \end{bmatrix}^{\!\top}
    + \log M
  \right)
  \begin{bmatrix} V_V \\ V_A \end{bmatrix},
  \label{eq:mot}
\end{equation}
with $M$ the visibility mask of \sref{sec:masks}. Each expert then continues
through its own post-attention path. The construction is the mixture of
transformers pattern \citep{liang2024mot} applied across modalities: the two
experts share the attention operator and keep separate parameters everywhere
else. Two structural conditions are required, namely $n_V = n_A$ and
$d_h^V = d_h^A$, so that the concatenated projections live in one space. The
hidden widths stay free, which lets a compact action expert couple to a wide
video backbone. Attention is parameter free, so the coupling costs exactly the parameters of the
second expert.

\Paragraph{Dual system with bridge cross-attention.}
The video expert runs to completion while hidden states at a set of bridge layers
$\mathcal{B} = \{\ell_1, \dots, \ell_J\}$ are collected. The action expert then
runs once, and its layer $j$ cross-attends to the bridge feature $h_V^{(\ell_j)}$,
\begin{equation}
  h_A^{(j)} \;\leftarrow\; h_A^{(j)} \;+\;
  \mathrm{CrossAttn}\!\left(h_A^{(j)},\; \mathrm{sg}^{\,\rho}\big[h_V^{(\ell_j)}\big]\right),
  \label{eq:bridge}
\end{equation}
where $\mathrm{sg}$ is the stop-gradient operator and $\rho \in \{0, 1\}$ selects
the detached variant. Setting $\rho = 1$ makes the video expert optimise the
video term of \eqref{eq:loss} alone, so its optimum is the pure video model and
the action expert becomes a read-only decoder of a fixed world representation.
Setting $\rho = 0$ lets the action term reshape the video representation, which
couples the two objectives through the shared trunk. The video expert reads its own stream alone in
both variants, which makes this family the one where video generation stays
exactly invariant to the action sample.

\Paragraph{Dual system with inverse-dynamics teacher forcing.}
Three branches pass through the joint attention in one forward: a noisy video
branch at level $\sigma_v$, a conditioning video branch at a second level
$\sigma_v^{\text{cond}}$, and a noisy action branch at level $\sigma_a$. The mask
blocks every cross-branch path except one,
\begin{equation}
  M^{\text{idm}} =
  \begin{bmatrix}
    M_{VV} & 0 & 0 \\
    0 & M_{VV} & 0 \\
    0 & \mathbf{1} & \mathbf{1}
  \end{bmatrix},
  \label{eq:idmmask}
\end{equation}
with rows and columns ordered as noisy video, conditioning video and action. Each
video branch is therefore self-contained and the action branch reads the
conditioning branch. The conditioning branch is noised with probability one half
so the action branch remains robust to an imperfect video prediction. Because the
video term is unconditional and the action term conditions on a video sample,
inference factors into two stages: denoise the video to completion, then denoise
the action against the finished video while reusing the video keys and values.

\Paragraph{Tri system.}
A third expert joins the joint attention. A frozen vision-language model
\citep{bai2025qwen25vl} encodes the observation and the instruction, a projector
maps its hidden states to a narrow understanding width, and the resulting tokens
enter the attention as a read-only tail that video and action queries may read.
The understanding expert supplies a semantic pathway beyond what the video
backbone's own text encoder carries, at the price of one narrow transformer plus
the projector.

\begin{figure}[t]
\centering
\adjustbox{max width=\linewidth}{%
\begin{tikzpicture}[font=\scriptsize,
  vb/.style={draw=tianzheblue!80!black, fill=tianzheblue!14, minimum width=9mm, minimum height=4.2mm, inner sep=0pt},
  ab/.style={draw=tianzheorange!80!black, fill=tianzheorange!22, minimum width=8mm, minimum height=4.2mm, inner sep=0pt},
  ub/.style={draw=tianzhepurple!80!black, fill=tianzhepurple!16, minimum width=8mm, minimum height=4.2mm, inner sep=0pt},
  cb/.style={draw=black!45, fill=black!7, minimum width=9mm, minimum height=4.2mm, inner sep=0pt},
  lk/.style={-{Stealth[length=1.3mm]}, draw=tianzhegreen!60!black},
  lk2/.style={{Stealth[length=1.3mm]}-{Stealth[length=1.3mm]}, draw=tianzhegreen!60!black},
  ttl/.style={font=\scriptsize\bfseries, align=center},
  cap/.style={font=\scriptsize, align=center, text=black!70}]

% ---------- panel A: single system ----------
\begin{scope}[shift={(0,0)}]
  \foreach \i in {0,1,2,3} { \node[vb] (a\i) at (0, \i*0.62) {}; }
  \node[ab, minimum width=9mm] (atok) at (0,-0.78) {};
  \node[cap] at (0,-1.60) {action and state\\tokens in sequence};
  \node[ttl] at (0,2.92) {(a) single system};
  \node[cap, anchor=west] at (0.62,1.86) {shared\\blocks};
  \draw[lk] (atok.north) -- (a0.south);
\end{scope}

% ---------- panel B: single system MoE ----------
\begin{scope}[shift={(4.2,0)}]
  \foreach \i in {0,1,2,3} { \node[vb] (b\i) at (0, \i*0.62) {}; }
  \foreach \i in {0,1,2,3} { \node[ab, minimum width=3.4mm] (e\i) at (0.92, \i*0.62) {}; }
  \foreach \i in {0,1,2,3} { \draw[lk2] (b\i.east) -- (e\i.west); }
  \node[ab, minimum width=9mm] (atok2) at (0,-0.78) {};
  \node[cap] at (0.3,-1.60) {expert feed-forward\\on the action segment};
  \node[ttl] at (0.4,2.92) {(b) single system\\with modality expert};
  \draw[lk] (atok2.north) -- (b0.south);
\end{scope}

% ---------- panel C: dual system joint self-attention ----------
\begin{scope}[shift={(8.6,0)}]
  \foreach \i in {0,1,2,3} { \node[vb] (c\i) at (0, \i*0.62) {}; }
  \foreach \i in {0,1,2,3} { \node[ab] (d\i) at (1.5, \i*0.62) {}; }
  \foreach \i in {0,1,2,3} { \draw[lk2] (c\i.east) -- (d\i.west); }
  \node[cap] at (0.75,-0.95) {one mixed attention\\at every layer};
  \node[ttl] at (0.75,2.92) {(c) dual system\\joint self-attention};
\end{scope}

% ---------- panel D: dual system bridge cross-attention ----------
\begin{scope}[shift={(0,-6.2)}]
  \foreach \i in {0,1,2,3} { \node[vb] (f\i) at (0, \i*0.62) {}; }
  \foreach \i in {0,1,2,3} { \node[ab] (g\i) at (1.5, \i*0.62) {}; }
  \foreach \i in {0,1,2,3} { \draw[lk] (f\i.east) -- (g\i.west); }
  \node[cap] at (0.75,-0.95) {video completes, then\\action reads the bridge};
  \node[ttl] at (0.75,2.92) {(d) dual system\\bridge cross-attention};
\end{scope}

% ---------- panel E: IDM teacher forcing ----------
\begin{scope}[shift={(4.6,-6.2)}]
  \foreach \i in {0,1,2,3} { \node[vb] (h\i) at (0, \i*0.62) {}; }
  \foreach \i in {0,1,2,3} { \node[cb] (k\i) at (1.35, \i*0.62) {}; }
  \foreach \i in {0,1,2,3} { \node[ab] (l\i) at (2.7, \i*0.62) {}; }
  \foreach \i in {0,1,2,3} { \draw[lk] (k\i.east) -- (l\i.west); }
  \node[cap] at (0,-0.95) {noisy};
  \node[cap] at (1.35,-0.95) {condition};
  \node[cap] at (2.7,-0.95) {action};
  \node[ttl] at (1.35,2.92) {(e) dual system\\teacher forcing};
\end{scope}

% ---------- panel F: tri system ----------
\begin{scope}[shift={(10.2,-6.2)}]
  \foreach \i in {0,1,2,3} { \node[vb] (m\i) at (0, \i*0.62) {}; }
  \foreach \i in {0,1,2,3} { \node[ab] (n\i) at (1.4, \i*0.62) {}; }
  \foreach \i in {0,1,2,3} { \node[ub] (o\i) at (2.8, \i*0.62) {}; }
  \foreach \i in {0,1,2,3} { \draw[lk2] (m\i.east) -- (n\i.west); }
  \foreach \i in {0,1,2,3} { \draw[lk] (o\i.west) -- (n\i.east); }
  \node[cap] at (1.4,-0.95) {frozen understanding\\tokens as a read-only tail};
  \node[ttl] at (1.4,2.92) {(f) tri system};
\end{scope}

% legend
\node[vb, label={[font=\scriptsize, label distance=-1mm]right:{\enspace video expert}}] at (0.1,-9.0) {};
\node[ab, label={[font=\scriptsize, label distance=-1mm]right:{\enspace action expert}}] at (3.4,-9.0) {};
\node[ub, label={[font=\scriptsize, label distance=-1mm]right:{\enspace understanding expert}}] at (7.0,-9.0) {};
\node[cb, label={[font=\scriptsize, label distance=-1mm]right:{\enspace conditioning branch}}] at (11.2,-9.0) {};
\end{tikzpicture}}
\caption{The six coupling architectures \sys{} exposes. Panels (a) and (b) place
action tokens inside the video sequence and reuse every video parameter. Panels
(c) to (e) give the action stream its own transformer and differ in where the two
streams meet. Panel (f) adds a frozen understanding expert. Double arrows mark a
symmetric meeting point and single arrows mark a one-directional read.}
\label{fig:families}
\end{figure}

\subsection{Cross-Modal Visibility as a Block Mask}
\label{sec:masks}

\Paragraph{The sequence and its blocks.}
Every coupling of \sref{sec:coupling} presents the attention with one sequence
\begin{equation}
  x \;=\; \big[\, \underbrace{x_V}_{S_V \text{ tokens}} \;\big|\;
                 \underbrace{x_A}_{S_A \text{ tokens}} \;\big|\;
                 \underbrace{x_R}_{S_R \text{ tokens}} \,\big],
  \label{eq:seq}
\end{equation}
where $x_R$ is a read-only tail carrying the state token of the single system or
the understanding tokens of the tri system. Visibility is a binary matrix
$M \in \{0,1\}^{S \times S}$ with $M_{ij} = 1$ when query $i$ reads key $j$, and
$\log M$ enters \eqref{eq:mot} as an additive bias. Writing $M$ in blocks,
\begin{equation}
  M \;=\;
  \begin{bmatrix}
    M_{VV} & M_{VA} & \mathbf{1} \\
    M_{AV} & \mathbf{1} & \mathbf{1} \\
    \mathbf{0} & \mathbf{0} & \mathbf{1}
  \end{bmatrix}.
  \label{eq:maskblocks}
\end{equation}
Three blocks are constant. Action tokens always read each other, since an action
chunk is a jointly denoised trajectory. The tail is readable by everything and
reads only itself, which keeps it a pure information source whose own
representation stays fixed. The two remaining blocks are the design space.

\Paragraph{Video self-visibility.}
$M_{VV}$ takes three values. The mode \texttt{bidirectional} sets every entry
to one. The mode \texttt{per-frame causal} builds a lower-triangular
matrix at frame granularity and expands it by the token count of one frame, so a
frame reads its own past. The mode \texttt{first-frame causal} sets
every entry to one and restricts the rows of the conditioning frame to
conditioning-frame keys, which keeps the clean prefix a self-contained
representation of the present observation. The reference configuration uses
\texttt{first-frame causal}.

\Paragraph{Cross-modal visibility.}
Let $P$ index the tokens of the clean conditioning frame. The four modes assign
the two cross blocks as follows.
\begin{equation}
  \begin{array}{llll}
    \textbf{mode} & \textbf{name} & M_{AV} & M_{VA} \\[2pt]
    \VAV  & \text{mutual}            & \mathbf{1}                & \mathbf{1} \text{ with rows } P \text{ zeroed} \\
    \AV   & \text{action reads video}& \mathbf{1}                & \mathbf{0} \\
    \VA   & \text{video reads action}& \text{columns } P \text{ only} & \mathbf{1} \text{ with rows } P \text{ zeroed} \\
    \ISO  & \text{isolated}          & \text{columns } P \text{ only} & \mathbf{0}
  \end{array}
  \label{eq:modes}
\end{equation}
Two invariants hold across all four modes. First, the rows of $P$ read only
video, so the keys and values of the conditioning prefix depend on the present
observation alone and stay valid for reuse across denoising iterations. Second,
$M_{AV}$ always grants the action stream access to the conditioning frame, so an
action is always grounded in the current observation even when the mode withholds
the predicted future. \fref{fig:masks} draws the resulting matrices.

\Paragraph{Why the best mode moves once pretraining enters.}
The two cross blocks differ in their consequences, because $M_{VA}$ decides
whether the video velocity field is a function of the action sample. That distinction becomes load bearing as soon as the pretraining mixture
contains footage whose action labels are absent.

\begin{proposition}[Consistency of an action-free mixture component]
\label{prop:actionfree}
Let the training mixture contain a component $\mathcal{D}_0$ on which the action
mask satisfies $M_A = \mathbf{0}$, so the action term of \eqref{eq:loss}
contributes no gradient. If $M_{VA} = \mathbf{0}$, then the video velocity field
$\hat{v}$ is measurable with respect to the video stream and the context alone,
and the video term evaluated on $\mathcal{D}_0$ estimates the same functional it
estimates on a labelled component. If $M_{VA} \neq \mathbf{0}$, then $\hat{v}$
depends on $a_{\sigma_a}$, whose conditional law on $\mathcal{D}_0$ is determined
by an unsupervised stream, and the video term on $\mathcal{D}_0$ estimates a
conditional functional whose conditioning variable carries no calibrated
information about the demonstrated motion.
\end{proposition}

\begin{proof}[Proof sketch]
With $M_{VA} = \mathbf{0}$ every video query assigns zero weight to every action
key at every layer, so the value written into a video residual stream is a
function of video and tail tokens only, and induction over layers gives
measurability of $\hat{v}$ with respect to that sub-sequence. The video term is
then the ordinary conditional flow-matching regression, whose target and
conditioning both come from $\mathcal{D}_0$ in the same form as from a labelled
component. With $M_{VA} \neq \mathbf{0}$ the same induction shows $\hat{v}$ is a
function of $a_{\sigma_a}$; on $\mathcal{D}_0$ the law of $a_0$ is unconstrained
by \eqref{eq:loss}, so the regression conditions on a variable whose relation to
the video target is set by the data source alone.
\end{proof}

Proposition~\ref{prop:actionfree} yields a concrete prediction. A model trained from
random initialisation on a mixture that contains action-free footage is best
served by a mode with $M_{VA} = \mathbf{0}$, which is \AV{} among the four. A
model whose action stream has already been calibrated by a large pretraining run
carries an action sample that is informative at every noise level, and restoring
$M_{VA}$ then supplies the video branch with a genuine additional conditioning
signal. The mutual mode \VAV{} therefore becomes the stronger configuration
after pretraining. \sref{sec:q4} measures both halves of this prediction.

\begin{figure}[t]
\centering
\adjustbox{max width=\linewidth}{%
\begin{tikzpicture}[font=\scriptsize, x=1mm, y=1mm,
  ttl/.style={font=\scriptsize\bfseries, align=center},
  lbl/.style={font=\tiny, text=black!70}]

\def\onst{fill=tianzhegreen!45, draw=black!35, line width=0.15pt}
\def\offst{fill=black!6, draw=black!25, line width=0.15pt}
\def\halfst{fill=tianzheblue!32, draw=black!35, line width=0.15pt}

% ---- top row: the four cross-modal modes ----
\def\maskmat#1#2#3#4#5#6#7{%
  \begin{scope}[shift={(#1,#2)}]
    \fill[#3] (0,10) rectangle (20,20);
    \fill[#4] (20,10) rectangle (30,20);
    \fill[#5] (0,0) rectangle (20,10);
    \fill[#6] (20,0) rectangle (30,10);
    \draw[black!60, line width=0.3pt] (0,0) rectangle (30,20);
    \draw[black!60, line width=0.3pt] (20,0) -- (20,20);
    \draw[black!60, line width=0.3pt] (0,10) -- (30,10);
    \node[ttl] at (15,26) {#7};
    \node[lbl, anchor=east] at (-1,15) {V};
    \node[lbl, anchor=east] at (-1,5) {A};
    \node[lbl] at (10,-3.5) {V};
    \node[lbl] at (25,-3.5) {A};
  \end{scope}}

\node[anchor=west, font=\scriptsize\itshape] at (0,36) {cross-modal visibility};
\maskmat{0}{0}{\onst}{\halfst}{\onst}{\onst}{$\VAV$\\mutual}
\maskmat{40}{0}{\onst}{\offst}{\onst}{\onst}{$\AV$\\action reads video}
\maskmat{80}{0}{\onst}{\halfst}{\halfst}{\onst}{$\VA$\\video reads action}
\maskmat{120}{0}{\onst}{\offst}{\halfst}{\onst}{$\ISO$\\isolated}

% ---- bottom row: the three video self-visibility modes ----
\node[anchor=west, font=\scriptsize\itshape] at (0,-12) {video self-visibility};

\begin{scope}[shift={(0,-44)}]
  \fill[\onst] (0,0) rectangle (18,18);
  \draw[black!60, line width=0.3pt] (0,0) rectangle (18,18);
  \node[ttl] at (9,23) {bidirectional};
\end{scope}

\begin{scope}[shift={(40,-44)}]
  \foreach \rr in {0,1,2} {
    \foreach \cc in {0,1,2} {
      \ifnum\cc>\rr
        \fill[\offst] (\cc*6,12-\rr*6) rectangle (\cc*6+6,18-\rr*6);
      \else
        \fill[\onst] (\cc*6,12-\rr*6) rectangle (\cc*6+6,18-\rr*6);
      \fi}}
  \draw[black!60, line width=0.3pt] (0,0) rectangle (18,18);
  \node[ttl] at (9,23) {per-frame causal};
\end{scope}

\begin{scope}[shift={(80,-44)}]
  \fill[\onst] (0,0) rectangle (18,18);
  \fill[\offst] (6,12) rectangle (18,18);
  \draw[black!60, line width=0.3pt] (0,0) rectangle (18,18);
  \node[ttl] at (9,23) {first-frame causal};
\end{scope}

\node[anchor=north west, align=left, font=\scriptsize] at (112,-24)
  {rows are queries, columns are keys\\[1pt]
   \textcolor{tianzhegreen!45}{$\blacksquare$}\enspace readable\\
   \textcolor{tianzheblue!32}{$\blacksquare$}\enspace readable outside the\\
   \hspace{3mm} conditioning rows\\
   \textcolor{black!12}{$\blacksquare$}\enspace withheld};
\end{tikzpicture}}
\caption{Visibility patterns. Top row: the four cross-modal modes of
\eqref{eq:modes} as block matrices over the video and action segments. Bottom
row: the three video self-visibility modes that fill $M_{VV}$. The conditioning
frame occupies the leading video rows and columns, and the mutual and
video-reads-action modes withhold the action block from exactly those rows, which
keeps the conditioning prefix reusable across denoising iterations.}
\label{fig:masks}
\end{figure}

\subsection{The Visual Representation the Model Denoises}
\label{sec:repr}

\Paragraph{The encoder interface.}
A visual encoder declares four structural properties: a latent channel width $C$,
a spatial compression factor $s$, a temporal compression factor, and whether the
latent admits a pixel decode. \sys{} fixes the temporal contract at
\begin{equation}
  T \;=\; 1 + \frac{N_v - 1}{4},
  \label{eq:temporal}
\end{equation}
so one latent frame carries the present observation and each further latent frame
summarises four future observation frames. With a patch size $(p_h, p_w)$ on the
backbone input projection, the video segment length is
\begin{equation}
  S_V \;=\; T \cdot \frac{H}{s\,p_h} \cdot \frac{W}{s\,p_w}.
  \label{eq:tokens}
\end{equation}
The reference configuration uses $N_v = 9$, $H = 384$, $W = 320$, $s = 16$ and
$(p_h, p_w) = (2,2)$, giving $T = 3$ and $S_V = 360$ video tokens against
$S_A = 32$ action tokens and one state token. Holding \eqref{eq:temporal} and
\eqref{eq:tokens} fixed across encoders is what makes the representation
comparison controlled, since every candidate then presents the backbone with the
same sequence length and the same temporal granularity.

\Paragraph{Two families of latent.}
Reconstructive latents come from a variational autoencoder trained to compress
pixels \citep{kingma2014vae,rombach2022ldm} and admit a pixel decode. Semantic
latents come from a self-supervised image or video encoder
\citep{simeoni2025dinov3,assran2025vjepa2}, are applied frame by frame, and are
pooled causally along time to satisfy \eqref{eq:temporal}. A parameter-free
per-token layer normalisation follows the pooling, which places the diffusion
target near an isotropic Gaussian and matches the assumption built into the
shifted noise grid. Semantic latents carry no decoder, so the backbone rebuilds
its input projection and its output head around $C$ and trains those two layers
from random initialisation.

\Paragraph{A reducer that separates width from semantics.}
The two families differ in channel width by more than an order of magnitude,
which confounds any direct comparison. \sys{} therefore inserts an optional
frozen reducer that compresses a semantic latent to the channel width of a
reconstructive latent. The reducer is a small per-token variational autoencoder
with a Gaussian bottleneck,
\begin{equation}
  q_\phi(w \mid y) \;=\; \mathcal{N}\!\big(\mu_\phi(y),\; \mathrm{diag}\,e^{\ell_\phi(y)}\big),
  \qquad \hat{y} \;=\; D_\phi(w),
  \label{eq:svaepost}
\end{equation}
trained offline on collected encoder features with the objective
\begin{equation}
  \mathcal{L}_{\text{red}}(\phi)
  \;=\;
  \big\| \hat{y} - y \big\|_2^2
  \;+\; \gamma \big(1 - \cos(\hat{y}, y)\big)
  \;+\; \beta \, \mathrm{KL}\!\big(q_\phi(w \mid y) \,\|\, \mathcal{N}(0, I)\big),
  \label{eq:svae}
\end{equation}
with $\gamma = 1$ and $\beta = 10^{-4}$, so reconstruction dominates and the
prior term only regularises the scale of the code. At evaluation the reducer
returns the posterior mean $\mu_\phi(y)$, which makes it a deterministic frozen
map. Setting the reducer output width to the reconstructive channel width leaves
the denoising problem identical in dimension, in temporal compression and in
sequence length, so the remaining difference between the two families is what the
channels encode. \sref{sec:q3} uses exactly this construction.

\subsection{One Observation Canvas and One Action Layout for Every Embodiment}
\label{sec:contract}

\Paragraph{Observation.}
Multiple camera views enter as a single composed canvas of height 384 and width
320. The head camera occupies a 256 by 320 tile at the top and two wrist cameras
occupy 128 by 160 tiles below it, and an absent view renders as a black tile.
One canvas per frame keeps the video stream single and keeps
\eqref{eq:tokens} independent of the camera count, so a platform with three
cameras and a platform with one camera present the backbone with the same
sequence.

\Paragraph{Action.}
All embodiments share one action width $d = 80$ with a fixed semantic slot
layout, drawn in \fref{fig:contract}. Slots 0 to 9 hold the left end-effector
position, its rotation in the six-dimensional continuous representation, and its
gripper aperture. Slots 10 to 33 hold twenty-four left hand joints. Slots 34 to
67 mirror both blocks for the right arm, and slots 68 to 79 stay reserved for
base motion and future extension. An embodiment $e$ with native action width
$d_e$ supplies an injective index map
$\pi_e : \{1, \dots, d_e\} \to \{1, \dots, 80\}$, and the framework applies the
scatter operator and its mask,
\begin{equation}
  (S_e y)_{\pi_e(i)} = y_i, \qquad
  (S_e y)_{j} = 0 \;\; \text{for } j \notin \mathrm{im}\,\pi_e,
  \qquad
  m_e[j] = \mathbbm{1}\{ j \in \mathrm{im}\,\pi_e \}.
  \label{eq:scatter}
\end{equation}
The action mask of \eqref{eq:loss} is then
$M_A = \mathbbm{1}_{H} m_e^{\!\top} \odot V$, with $V$ the horizon validity mask
that removes padded steps.

\begin{proposition}[Slot isolation]
\label{prop:slots}
Consider a batch containing samples from embodiments $e_1, \dots, e_n$ and let
$\theta_j$ collect the parameters of the output head that produce slot $j$. The
gradient of \eqref{eq:loss} with respect to $\theta_j$ is a sum over the samples
whose mask satisfies $m_{e}[j] = 1$. Slots that no embodiment in the batch maps
receive zero gradient, and the contribution of each embodiment equals its native
per-dimension loss up to the per-sample normalisation $1 / \sum M_A$.
\end{proposition}

\begin{proof}[Proof sketch]
$M_A$ enters \eqref{eq:loss} as an elementwise factor, so the residual at slot
$j$ of a sample with $m_e[j] = 0$ is multiplied by zero before the sum. The
derivative of the mask-normalised norm with respect to $\theta_j$ therefore
retains only the terms with $m_e[j] = 1$, and within those terms the summand is
the native squared residual divided by the count of valid entries.
\end{proof}

Proposition~\ref{prop:slots} is what allows a mixture of gripper platforms, dexterous hands
and single-arm setups to train one action head at one width. Normalisation
completes the contract. Position and joint dimensions are rescaled to
$[-1, 1]$ by the per-source minimum and maximum, rotation blocks are pinned to
the identity range so the continuous rotation representation passes through
unchanged, and gripper aperture is canonicalised so that the low value denotes a
closed gripper and the high value denotes an open one. Every downstream platform
inverts this map at deployment.

\begin{figure}[t]
\centering
\adjustbox{max width=\linewidth}{%
\begin{tikzpicture}[font=\scriptsize, x=1mm, y=1mm,
  sl/.style={draw=black!55, line width=0.3pt},
  tl/.style={draw=black!55, line width=0.3pt},
  lbl/.style={font=\tiny, align=center}]

% ---- action slot layout ----
\node[anchor=west, font=\scriptsize\itshape] at (0,20) {unified action layout, width 80};
\filldraw[fill=tianzheblue!35, draw=black!55, line width=0.3pt] (0,6) rectangle (12,14);
\filldraw[fill=tianzheblue!18, draw=black!55, line width=0.3pt] (12,6) rectangle (44,14);
\filldraw[fill=tianzheorange!35, draw=black!55, line width=0.3pt] (44,6) rectangle (56,14);
\filldraw[fill=tianzheorange!18, draw=black!55, line width=0.3pt] (56,6) rectangle (88,14);
\filldraw[fill=black!8, draw=black!55, line width=0.3pt] (88,6) rectangle (104,14);
\draw[sl] (0,6) rectangle (104,14);

\node[lbl] at (6,10) {left arm\\pose and grip};
\node[lbl] at (28,10) {left hand joints};
\node[lbl] at (50,10) {right arm\\pose and grip};
\node[lbl] at (72,10) {right hand joints};
\node[lbl] at (96,10) {reserved};

\foreach \x/\t in {0/0, 12/10, 44/34, 56/44, 88/68, 104/80} {
  \draw[black!55, line width=0.3pt] (\x,6) -- (\x,3.6);
  \node[font=\tiny, anchor=north] at (\x,3.4) {\t};}

% ---- observation canvas ----
\begin{scope}[shift={(118,-8)}]
  \node[anchor=west, font=\scriptsize\itshape] at (0,28) {observation canvas};
  \filldraw[fill=tianzhegreen!22, draw=black!55, line width=0.3pt] (0,10) rectangle (32,35.6);
  \filldraw[fill=tianzhegreen!12, draw=black!55, line width=0.3pt] (0,-2.8) rectangle (16,10);
  \filldraw[fill=black!8, draw=black!55, line width=0.3pt] (16,-2.8) rectangle (32,10);
  \node[lbl] at (16,22.8) {head camera\\256 $\times$ 320};
  \node[lbl] at (8,3.6) {left wrist\\128 $\times$ 160};
  \node[lbl] at (24,3.6) {absent view\\rendered black};
  \draw[tl] (0,-2.8) rectangle (32,35.6);
  \draw[{Stealth[length=1.4mm]}-{Stealth[length=1.4mm]}, black!60] (-3.4,-2.8) -- (-3.4,35.6)
    node[midway, anchor=east, font=\tiny, rotate=90, yshift=2mm] {384};
  \draw[{Stealth[length=1.4mm]}-{Stealth[length=1.4mm]}, black!60] (0,-6.2) -- (32,-6.2)
    node[midway, anchor=north, font=\tiny] {320};
\end{scope}
\end{tikzpicture}}
\caption{The cross-embodiment contract. Left: the fixed action layout of width
80, with tick marks giving slot boundaries. An embodiment writes into the slots
its hardware measures and the per-dimension mask keeps every other slot out of
the loss. Right: the composed observation canvas that carries up to three camera
views inside one video stream at a fixed token count.}
\label{fig:contract}
\end{figure}

\subsection{Scheduling the Two Streams at Deployment}
\label{sec:deploy}

Proposition~\ref{prop:product} established that any monotone path through the level square
is admissible. This subsection parameterises a family of such paths, then
specifies how a generated chunk is executed against a running control loop.

\Paragraph{A two-parameter family of denoising trajectories.}
Let $u_k = k / K'$ measure global progress across $K'$ iterations and let
cleanness denote the complement of a stream's noise level before shifting. One
stream leads and the other lags. The lead stream takes cleanness
$c_{\text{lead}}(u) = f_\alpha(u)$ with $\alpha \geq 1$, and the lag stream takes
cleanness $c_{\text{lag}}(u) = \mathrm{clip}\big((u - \delta) / (1 - \delta), 0, 1\big)$
with $0 \leq \delta < 1$. Each stream's noise level then rides its own training
grid,
\begin{equation}
  \sigma_{\text{lead}}(u) = f_{\lambda}\!\big(1 - f_\alpha(u)\big),
  \qquad
  \sigma_{\text{lag}}(u) = f_{\lambda}\!\big(1 - c_{\text{lag}}(u)\big).
  \label{eq:traj}
\end{equation}

\begin{lemma}[Properties of the shift map]
\label{lem:shift}
For $\alpha > 0$ the map $f_\alpha$ of \eqref{eq:shift} is a monotone increasing
bijection of $[0,1]$ with $f_\alpha(0) = 0$ and $f_\alpha(1) = 1$; its inverse is
$f_{1/\alpha}$; and $1 - f_\alpha(1 - t) = f_{1/\alpha}(t)$ for all $t$.
Moreover $f_\alpha(u) \geq u$ for every $\alpha \geq 1$.
\end{lemma}

\begin{proof}
The endpoint values follow by substitution. Differentiating gives
$f_\alpha'(t) = \alpha / (1 + (\alpha-1)t)^2 > 0$, which gives monotonicity, and
solving $w = \alpha t / (1 + (\alpha - 1)t)$ for $t$ gives
$t = w / (\alpha - (\alpha - 1) w) = f_{1/\alpha}(w)$. For the reflection
identity,
\[
  1 - f_\alpha(1-t)
  = \frac{1 + (\alpha-1)(1-t) - \alpha(1-t)}{1 + (\alpha-1)(1-t)}
  = \frac{t}{\alpha - (\alpha-1)t}
  = f_{1/\alpha}(t).
\]
Finally
$f_\alpha(u) - u = u(1-u)(\alpha-1) / \big(1 + (\alpha-1)u\big) \geq 0$
whenever $\alpha \geq 1$ and $u \in [0,1]$.
\end{proof}

Lemma~\ref{lem:shift} gives the family three properties that matter in practice. The
lead stream reaches any level of cleanness at a smaller global progress than the
lag stream, so it finishes earlier. The reflection identity lets
\eqref{eq:traj} be computed as $f_{1/\alpha}$ applied to the base grid, so at
$\alpha = 1$ and $\delta = 0$ the two curves coincide with the base grid and the
trajectory reduces exactly to the synchronous diagonal. Both curves stay
monotone and terminate at zero, so \eqref{eq:euler} remains a valid Euler
integration of the learned field. Selecting the lead stream, $\alpha$ and
$\delta$ therefore spans a continuum whose endpoints are a video-leading
trajectory and an action-leading trajectory, and whose centre is the synchronous
one. \fref{fig:sched} draws the resulting paths.

\Paragraph{Executing a chunk against a control loop.}
A generation returns $H$ actions and the executor consumes $h \leq H$ of them
before regenerating, which is receding-horizon control with a replanning period
of $h$ steps. The synchronous executor blocks on the generation, so its control
period alternates between one generation latency and $h-1$ cheap steps. The
asynchronous executor keeps a single background generation in flight: when
$h - d$ actions remain in the buffer it launches the next generation, where $d$
is the expected latency measured in action steps. A generation launched at
control step $t_s$ and adopted at step $t$ has its first $t - t_s$ actions
already superseded, so the executor discards them and enqueues the remaining
window. Writing $\Delta t$ for the control period and $T_{\text{gen}}$ for the
generation latency, execution stays free of stalls whenever
\begin{equation}
  T_{\text{gen}} \;\leq\; d \cdot \Delta t
  \qquad\text{and}\qquad
  d \;<\; h \;\leq\; H,
  \label{eq:async}
\end{equation}
and the number of actions usable from an adopted chunk is $H - (t - t_s)$.

\Paragraph{Reducing the cost of a generation.}
Three deployment optimisations act inside Algorithm~\ref{alg:sample}. The velocity cache
compares consecutive video velocity predictions and reuses the previous
prediction when their cosine similarity exceeds a threshold $\tau$, subject to a
budget of $m$ consecutive reuses, so a sampler of $K'$ iterations performs at
most $K'$ and at least $\lceil K' / (m+1) \rceil$ video forwards. Prompt
embeddings are cached across requests, since a benchmark episode repeats one
instruction. Fixed-shape compilation is applied to the per-layer joint loop of
each coupling family, which removes per-layer dispatch overhead at the price of a
first-request warm-up. Skipping the pixel decode removes the decoder from the
critical path entirely, which is admissible because the control loop consumes
only the action stream.

\begin{figure}[t]
\centering
\adjustbox{max width=\linewidth}{%
\begin{tikzpicture}[font=\scriptsize]

% ---------- (a) the level square ----------
\begin{scope}
  \draw[->, black!70] (0,0) -- (5.3,0) node[anchor=north east, font=\scriptsize] {$\sigma_v$};
  \draw[->, black!70] (0,0) -- (0,5.3) node[anchor=south west, font=\scriptsize] {$\sigma_a$};
  \foreach \t/\l in {0/0, 5/1} {
    \node[anchor=north, font=\tiny] at (\t,-0.05) {\l};
    \node[anchor=east, font=\tiny] at (-0.05,\t) {\l};}
  \draw[black!18, line width=0.3pt] (5,0) -- (5,5) -- (0,5);

  % synchronous diagonal
  \draw[tianzhegreen!60!black, line width=1.1pt] (5,5) -- (0,0);
  % video leading: sigma_v drops faster
  \draw[tianzheblue!75!black, line width=1.1pt, densely dashed]
    (5,5) .. controls (2.2,4.7) and (0.6,3.2) .. (0,0);
  % action leading: sigma_a drops faster
  \draw[tianzheorange!70!black, line width=1.1pt, densely dotted]
    (5,5) .. controls (4.7,2.2) and (3.2,0.6) .. (0,0);

  \node[anchor=west, font=\tiny, text=tianzhegreen!50!black] at (2.55,2.15) {synchronous};
  \node[anchor=west, font=\tiny, text=tianzheblue!60!black] at (0.35,3.95) {video leading};
  \node[anchor=west, font=\tiny, text=tianzheorange!60!black] at (2.95,0.75) {action leading};
  \node[anchor=north, font=\scriptsize] at (2.65,-0.75) {(a) trajectories through the level square};
  \node[circle, fill=black!70, inner sep=0.7pt] at (5,5) {};
  \node[circle, fill=black!70, inner sep=0.7pt] at (0,0) {};
  \node[anchor=south west, font=\tiny] at (4.55,4.9) {start};
  \node[anchor=south west, font=\tiny] at (0.1,0.05) {clean};
\end{scope}

% ---------- (b) execution timelines ----------
\begin{scope}[shift={(8.2,0)}]
  \node[anchor=north, font=\scriptsize] at (3.6,-0.75) {(b) synchronous and asynchronous execution};

  % synchronous
  \node[anchor=east, font=\tiny] at (-0.15,4.15) {synchronous};
  \filldraw[fill=tianzhered!30, draw=black!45, line width=0.3pt] (0,3.75) rectangle (1.5,4.55);
  \node[font=\tiny] at (0.75,4.15) {generate};
  \foreach \k in {0,...,5} {
    \filldraw[fill=tianzhegreen!30, draw=black!45, line width=0.3pt] (1.5+\k*0.4,3.75) rectangle (1.9+\k*0.4,4.55);}
  \filldraw[fill=tianzhered!30, draw=black!45, line width=0.3pt] (3.9,3.75) rectangle (5.4,4.55);
  \node[font=\tiny] at (4.65,4.15) {generate};
  \foreach \k in {0,...,3} {
    \filldraw[fill=tianzhegreen!30, draw=black!45, line width=0.3pt] (5.4+\k*0.4,3.75) rectangle (5.8+\k*0.4,4.55);}
  \node[anchor=west, font=\tiny, text=tianzhered!70!black] at (7.15,4.15) {loop stalls};

  % asynchronous
  \node[anchor=east, font=\tiny] at (-0.15,2.15) {asynchronous};
  \foreach \k in {0,...,17} {
    \filldraw[fill=tianzhegreen!30, draw=black!45, line width=0.3pt] (\k*0.4,1.75) rectangle (0.4+\k*0.4,2.55);}
  \filldraw[fill=tianzheblue!25, draw=black!45, line width=0.3pt, rounded corners=1pt] (1.6,0.75) rectangle (3.1,1.45);
  \node[font=\tiny] at (2.35,1.1) {background generate};
  \filldraw[fill=tianzheblue!25, draw=black!45, line width=0.3pt, rounded corners=1pt] (5.2,0.75) rectangle (6.7,1.45);
  \node[font=\tiny] at (5.95,1.1) {background generate};
  \draw[-{Stealth[length=1.3mm]}, black!60] (2.35,1.45) -- (2.35,1.72);
  \draw[-{Stealth[length=1.3mm]}, black!60] (5.95,1.45) -- (5.95,1.72);
  \draw[-{Stealth[length=1.3mm]}, black!60] (1.6,1.72) -- (1.6,1.48);
  \draw[-{Stealth[length=1.3mm]}, black!60] (5.2,1.72) -- (5.2,1.48);
  \node[anchor=west, font=\tiny, text=tianzheblue!60!black] at (7.15,1.1) {launch at $h-d$};
  \node[anchor=west, font=\tiny, text=tianzhegreen!50!black] at (7.15,2.15) {uniform control period};
\end{scope}
\end{tikzpicture}}
\caption{Deployment schedules. (a) The three trajectories of \eqref{eq:traj}
inside the level square, with $\alpha = 1$ recovering the synchronous diagonal
exactly by Lemma~\ref{lem:shift}. (b) The synchronous executor blocks the control loop
for one generation latency per replanning period, and the asynchronous executor
launches the next generation while $d$ actions remain, which keeps the control
period uniform under condition \eqref{eq:async}.}
\label{fig:sched}
\end{figure}
